\documentclass[aspectratio=169]{beamer}

\usepackage{xeCJK}
\setCJKmainfont[BoldFont=FandolHei-Regular.otf]{FandolSong-Regular.otf}
\setCJKsansfont[BoldFont=FandolHei-Regular.otf]{FandolHei-Regular.otf}
\setCJKmonofont[BoldFont=FandolHei-Regular.otf]{FandolFang-Regular.otf}

\usetheme{USST}
\usepackage{amsmath,amssymb,booktabs,array,graphicx}

\newcommand{\titleplaceholder}{上海理工大学汇报模板}
\newcommand{\subtitleplaceholder}{适用于课程展示、读书报告与学术汇报}
\newcommand{\authorplaceholder}{姓名}
\newcommand{\instituteplaceholder}{学院}

\title[\titleplaceholder]{\titleplaceholder}
\subtitle{\subtitleplaceholder}
\author{\authorplaceholder}
\institute{\instituteplaceholder}
\date{\today}

\begin{document}

\ussttitleframe
\usstoutlineframe

\section{模板说明}
\begin{frame}{模板说明}
\large
这是一套为上海理工大学学生整理的 Beamer 汇报模板，保留了校徽与主视觉风格，也尽量让排版保持简洁、正式和清晰。

\vspace{0.3cm}
\begin{itemize}
  \item 修改封面信息即可快速开始使用
  \item 使用 \texttt{\textbackslash section\{\}} 可自动生成目录
  \item 建议每页只表达一个核心信息
\end{itemize}
\end{frame}

\section{版式示例}
\begin{frame}{双栏版式示例}
\begin{columns}[T,onlytextwidth]
\column{0.56\textwidth}
\large
\textbf{左侧内容区}
\begin{itemize}
  \item 适合放研究背景、方法步骤或结论概述
  \item 文字建议分成 3 到 5 个短点
  \item 段落过长时可以拆成两页展示
\end{itemize}

\column{0.40\textwidth}
\large
\textbf{右侧内容区}
\begin{itemize}
  \item 适合放图片、对比要点或补充说明
  \item 也可以替换成图表、流程图或公式
\end{itemize}
\end{columns}
\end{frame}

\begin{frame}{公式与定义示例}
\large
这页适合展示定义、公式或关键结论：
\[
f(x)=x^2+2x+1,\qquad x^\star=-1
\]

\vspace{0.25cm}
\begin{itemize}
  \item 公式下方建议补一句解释，帮助听众快速理解
  \item 如果公式较长，可以只保留最关键的一步
\end{itemize}
\end{frame}

\section{内容建议}
\begin{frame}{表格示例}
\small
\begin{tabular}{>{\raggedright\arraybackslash}p{0.28\textwidth}>{\raggedright\arraybackslash}p{0.28\textwidth}>{\raggedright\arraybackslash}p{0.28\textwidth}}
\toprule
\textbf{页面类型} & \textbf{适合内容} & \textbf{使用建议}\\
\midrule
封面页 & 题目、姓名、单位 & 信息保持简洁\\
目录页 & 汇报结构 & 控制在 4 到 6 部分\\
总结页 & 结论与展望 & 用短句强化重点\\
\bottomrule
\end{tabular}

\vspace{0.35cm}
\large
如果你的内容以图片或实验结果为主，也可以把表格页直接换成图文页。
\end{frame}

\section{结束总结}
\begin{frame}{总结}
\large
\begin{itemize}
  \item 这套模板保留了上海理工大学的视觉识别元素
  \item 结构清晰，适合日常课程与学术汇报场景
  \item 你可以在此基础上继续替换文字、图片与配色细节
\end{itemize}
\end{frame}

\usstthanksframe

\end{document}
